\subsection{代数 A 到 A-模的导子}

我们在本号中假设第 2 号的情形 (II) 成立。那么存在一个分次 K-代数 A 和一个分次 K-模 E，以及两个次数为 0 的 K-线性映射

\[
\mathbf {E} \otimes_ {\mathbf {K}} \mathbf {A} \rightarrow \mathbf {E}, \quad \mathbf {A} \otimes_ {\mathbf {K}} \mathbf {E} \rightarrow \mathbf {E}
\]

记作

\[
x \otimes a \mapsto x. a \quad \text { and } \quad a \otimes x \mapsto a. x \quad \text { for   } a \in A \text {   and   } x \in E.
\]

命题 3. 设 \(d:A\to E\) 是一个 \(\varepsilon\) -导子（次数为 \(\delta\) ）。则 \(\operatorname{Ker}(d)\) 是 A 的一个分次子代数；若 A 有单位元，则它属于 \(\operatorname{Ker}(d)\) 。

显然 \(\operatorname{Ker}(d)\) 是 A 的一个分次子 K-模；此外，第 2 号的关系 (8) 表明，若 x 和 y 是属于 \(\operatorname{Ker}(d)\) 的两个齐次元素，则 \(d(xy) = 0\) ，从而 \(xy \in \operatorname{Ker}(d)\) 。最后，若 A 有单位元 1（次数为 0，参见 § 3 第 1 号），则把第 2 号的关系 (8) 中的 x 和 y 都换成 1，就得到 \(d(1) = d(1) + d(1)\) ，从而 \(d(1) = 0\) 。

推论。设 \(d_1\) 和 \(d_2\) 是两个 \(\varepsilon\) -导子（从 \(A\) 到 \(E\) ，次数同为 \(\delta\) ）。若 \(d_1\) 和 \(d_2\) 在代数 \(A\) 的一个生成系的每个元素上取相同的值，则 \(d_1 = d_2\) 。

\(d_{1} - d_{2}\) 是一个 \(\varepsilon\) -导子（次数为 \(\delta\) ），因此 \(\operatorname{Ker}(d_1 - d_2)\) 是 A 的一个子代数，它包含 A 的一个生成系，从而等于 A。

命题 4. 设 \(d:A\to E\) 是一个 \(\varepsilon\) -导子（次数为 \(\delta\) ）。假设 A 有单位元 1，并设 \(x\) 是 A 的一个齐次元素，它在 A 中有逆元 \(x^{-1}\) 。则

\[
d (x ^ {- 1}) = - \varepsilon_ {\delta , \deg (x)} x ^ {- 1} ((d x) x ^ {- 1}) = - \varepsilon_ {\delta , \deg (x)} (x ^ {- 1} (d x)) x ^ {- 1}.\tag{19}
\]

我们有 \(d(xx^{-1}) = d(1) = 0\) （命题 3），由此得出

\[
(d x) x ^ {- 1} + \varepsilon_ {\delta , \deg (x)} x (d (x ^ {- 1})) = 0
\]

这证明了 (19) 中的第一个公式。另一方面， \(x^{-1}\) 是次数为 \(-\deg(x)\) 的齐次元素，且由第 1 号的公式 (1) 有 \(\varepsilon_{\delta,\deg(x)} = \varepsilon_{\delta,-\deg(x)}\) ；写出 \(d(x^{-1}x) = 0\) ，即类似地得到 (19) 的第二个公式。

命题 5. 设 A 是一个整环，并令 L 为其分式域。A 到 L 上向量空间 E（视为 A-模）的每个导子都可以唯一地扩张为 L 到 E 的导子。

设 \(d\) 是 \(\mathbf{A}\) 到 \(\mathbf{E}\) 的一个导子， \(\bar{d}\) 是 \(\mathbf{L}\) 到 \(\mathbf{E}\) 的一个导子且扩张 \(d\) ；那么，对 \(u \in \mathbf{A}, v \in \mathbf{A}, v \neq 0\) ，必然有，根据 (19)，

\[
\bar {d} (u / v) = v ^ {- 1} d u - u v ^ {- 2} d v\tag{20}
\]

这就证明了 \(\bar{d}\) 的唯一性。反之，我们证明 \(\bar{d}\) 可以由公式 (20) 定义；首先必须验证，若 \(u / v = u' / v'\) ，则当把 \(u\) 换成 \(u'\) 、把 \(v\) 换成 \(v'\) 时，(20) 右端的值不改变。现在， \(uv' = vu'\) ，因此 \(v'(du) + u(dv') = v(du') + u'(dv)\) ，从而 \(v'(du - uv^{-1}dv) = v(du' - u'v'^{-1}dv')\) ，因为 \(uv'v^{-1} = u'\) 且 \(u'v'^{-1}v = u\) 。于是定义了一个映射 \(\bar{d}: L \to E\) ，它扩张 \(d\) ；立即可验证它是 K-线性的且是一个导子。

命题 6. 设 A 是一个含单位元的结合分次 K-代数，E 是一个分次 (A, A)-双模。若 \(d: \mathbf{A} \to \mathbf{E}\) 是一个 \(\varepsilon\) -导子（次数为 \(\delta\) ），则对 A 的齐次元素的每个有限序列 \((x_i)_{1 \leqslant i \leqslant n}\) （次数分别为 \(\xi_i\) ， \(1 \leqslant i \leqslant n\) ），

\[
d \left(x _ {1} x _ {2} \dots x _ {n}\right) = \sum_ {i = 1} ^ {n} \varepsilon_ {\delta, \xi_ {1} + \dots + \xi_ {i - 1}} x _ {1} \dots x _ {i - 1} \left(d x _ {i}\right) x _ {i + 1} \dots x _ {n}.\tag{21}
\]

公式 (21) 对 \(n = 0\) 是平凡的，并对 \(n\) 用归纳法证明，其中用到 (4)（第 2 号）。

推论。设 A 是一个含幺交换结合代数，E 是一个 A-模。若 \(d: \mathbf{A} \to \mathbf{E}\) 是一个导子，则对于每个整数 \(n \geqslant 0\) ,

\[
d (x ^ {n}) = n x ^ {n - 1} (d x) \quad \text { for all } x \in \mathbf {A}.\tag{22}
\]

只需给 A 赋予平凡分次，并在 (21) 中取所有 \(x_{i}\) 等于 \(x\) 即可。

我们回到一个 \(\varepsilon\) -导子 \(d: \mathbf{A} \to \mathbf{E}\) （次数为 \(\delta\) ）的一般情形。设 \(Z_{\varepsilon}\) 是这样的 \(a \in \mathbf{A}\) 组成的集合：使得对每个齐次分量 \(a_{\alpha}\) （属于 \(a\) ，次数为 \(\alpha\) ），对每个齐次的 \(x\) （在 \(\mathbf{E}\) 中），

\[
x a _ {\alpha} = \varepsilon_ {\alpha , \deg (x)} a _ {\alpha} x.\tag{23}
\]

若 A 是一个含单位元的结合分次代数，且 E 是一个分次 (A, A)-双模，则由该定义立即得出 \(Z_{\varepsilon}\) 是A的一个分次子代数，且包含单位元。

命题 7. 设 A 是一个含单位元的结合分次代数，E 是一个分次 (A, A)-双模。令 \(d: \mathbf{A} \to \mathbf{E}\) 是一个 \(\varepsilon\) -导子（次数为 \(\delta\) ），并令 \(a\) 是 \(Z_{\varepsilon}\) 的一个次数为 \(\alpha\) 的齐次元素。则映射 \(x \mapsto a(dx)\) 是一个 \(\varepsilon\) -导子（次数为 \(\delta + \alpha\) ）。

我们记 \(d'(x) = a(dx)\) ；对齐次元素 \(x\) （次数为 \(\xi\) ， \(A\) 中的）以及 \(y \in A\) ，根据 (23) 和 (1)（第 1 号），

\[
\begin{array}{r l} d ^ {\prime} (x y) & = a ((d x) y) + \varepsilon_ {\delta , \xi} a (x (d y)) = (a (d x)) y + \varepsilon_ {\delta + \alpha , \xi} (x a) (d y) \\ & = (d ^ {\prime} x) y + \varepsilon_ {\delta + \alpha , \xi} x (d ^ {\prime} y). \end{array}
\]

命题 7 表明，A 到 E 的 \(\varepsilon\) -导子组成的 K-模是一个分次 \(Z_{\varepsilon}\) -模（类型为 \(\Delta\) ）。

